\documentclass[aspectratio=169,10pt]{beamer}
\makeatletter\def\input@path{{../../estilo/}{../estilo/}{estilo/}}\makeatother
\usepackage{devmobile}
\modulo[Kotlin]{15}{Publicação na Google Play}{Build de release, assinatura, App Bundle, Play Console, faixas de teste, verificação de desenvolvedor, LGPD e CI/CD}

\begin{document}
\capa

\begin{frame}{Objetivos do módulo}
  \begin{conceito}[Ao final deste módulo você será capaz de]
    \begin{itemize}
      \item Preparar um build de release otimizado e versionado.
      \item Assinar o app e entender a Play App Signing.
      \item Criar a conta e a ficha do app na Play Console, incluindo as declarações obrigatórias.
      \item Conduzir testes interno, fechado e aberto e cumprir a regra de 12 testadores.
      \item Lançar de forma gradual e acompanhar a saúde do app em produção.
      \item Atender à LGPD e automatizar a publicação com CI/CD.
    \end{itemize}
  \end{conceito}
\end{frame}

\sumario

% =============================================================================
%  Fundamentos compartilhados --- Módulo 15: distribuição de apps
%  Incluído pelas trilhas Kotlin e Flutter com % =============================================================================
%  Fundamentos compartilhados --- Módulo 15: distribuição de apps
%  Incluído pelas trilhas Kotlin e Flutter com % =============================================================================
%  Fundamentos compartilhados --- Módulo 15: distribuição de apps
%  Incluído pelas trilhas Kotlin e Flutter com \input{../fundamentos/f15-publicacao}
% =============================================================================
\section{Fundamentos da distribuição de apps}

\begin{frame}{Do código ao aparelho do usuário}
  \centering
  \begin{tikzpicture}[node distance=2.4mm,every node/.style={font=\scriptsize,minimum height=9mm,minimum width=1.62cm}]
    \node[caixa] (c) {Código\\ \tiny versionado};
    \node[caixapetroleo,right=of c] (b) {Build de\\ release};
    \node[caixapetroleo,right=of b] (a) {Assinatura};
    \node[caixaambar,right=of a] (e) {Envio\\ à loja};
    \node[caixaambar,right=of e] (r) {Revisão\\ e políticas};
    \node[caixa,right=of r] (t) {Faixas\\ de teste};
    \node[caixaescura,right=of t] (p) {Produção};
    \foreach \x/\y in {c/b,b/a,a/e,e/r,r/t,t/p} \draw[seta] (\x) -- (\y);
    \draw[setaambar,rounded corners] (p.south) -- ++(0,-6mm) -| node[rotulo,pos=0.25,below]{\tiny monitorar falhas e avaliações $\rightarrow$ nova versão} (c.south);
  \end{tikzpicture}
  \vspace{6pt}
  \begin{columns}[T]
    \begin{column}{0.48\textwidth}
      \scriptsize
      \begin{tabularx}{\linewidth}{@{}l L L@{}}
        \toprule
        & \cabecalho{Debug} & \cabecalho{Release} \\
        \midrule
        Objetivo & desenvolver & distribuir \\
        Código & com verificações e depuração & otimizado, reduzido e ofuscável \\
        Desempenho & mais lento & o real \\
        Assinatura & chave de depuração automática & sua chave de release \\
        \bottomrule
      \end{tabularx}
    \end{column}
    \begin{column}{0.48\textwidth}
      \begin{definicao}[Build de release]
        Versão do app compilada para o usuário final: sem ferramentas de depuração, com código otimizado e assinada com a chave do desenvolvedor.
      \end{definicao}
      \begin{atencao}
        Teste sempre o build de release antes de publicar: ofuscação e otimizações podem expor defeitos que não aparecem em debug.
      \end{atencao}
    \end{column}
  \end{columns}
\end{frame}

\begin{frame}{Assinatura digital: a teoria}
  \begin{columns}[c]
    \begin{column}{0.54\textwidth}
      \centering
      \begin{tikzpicture}[node distance=4mm,every node/.style={font=\scriptsize}]
        \node[caixa,minimum width=1.6cm] (app) {App};
        \node[caixa,right=of app,minimum width=1.6cm] (h) {Hash\\ \tiny resumo único};
        \node[caixaescura,right=of h,minimum width=2cm] (s) {Assinatura};
        \node[caixaambar,above=5mm of s,minimum width=2cm] (pk) {Chave \textbf{privada}\\ \tiny só o dev tem};
        \draw[seta] (app) -- (h); \draw[seta] (h) -- (s); \draw[seta] (pk) -- (s);
        \node[caixa,below=9mm of app,minimum width=1.6cm] (app2) {App\\ \tiny recebido};
        \node[caixa,right=of app2,minimum width=1.6cm] (h2) {Hash\\ \tiny recalculado};
        \node[caixapetroleo,right=of h2,minimum width=2cm] (v) {Confere?};
        \node[caixa,below=5mm of v,minimum width=2cm] (pub) {Chave \textbf{pública}\\ \tiny no certificado};
        \draw[seta] (app2) -- (h2); \draw[seta] (h2) -- (v); \draw[seta] (pub) -- (v);
        \draw[setaambar] (s.south) -- node[rotulo,right]{\tiny vai junto} (v.north);
        \node[rotulo,left=1mm of app,font=\tiny,rotate=90,anchor=south] {quem publica};
        \node[rotulo,left=1mm of app2,font=\tiny,rotate=90,anchor=south] {sistema};
      \end{tikzpicture}
    \end{column}
    \begin{column}{0.44\textwidth}
      \begin{definicao}[Assinatura digital]
        Resumo (\emph{hash}) do pacote cifrado com a \textbf{chave privada} do desenvolvedor. Qualquer um verifica com a \textbf{chave pública}, mas só quem tem a privada consegue assinar.
      \end{definicao}
      \begin{conceito}[O que ela garante]
        \textbf{Integridade}: o pacote não foi alterado.\\
        \textbf{Autenticidade}: a atualização vem do mesmo autor do app instalado.
      \end{conceito}
      \begin{analogia}
        O lacre de cera com o sinete do remetente: qualquer um reconhece o selo, ninguém o falsifica sem o anel.
      \end{analogia}
    \end{column}
  \end{columns}
\end{frame}

\begin{frame}{Assinatura em cada plataforma}
  \begin{columns}[T]
    \begin{column}{0.5\textwidth}
      \begin{conceito}[Android]
        O pacote é assinado com uma chave guardada num \emph{keystore}. Com o \textbf{Play App Signing}, você assina com a \textbf{chave de upload}; o Google confere e reassina com a \textbf{chave do app}, que fica guardada com ele.
      \end{conceito}
      \centering
      \begin{tikzpicture}[node distance=9mm,every node/.style={font=\tiny,align=center}]
        \node[caixa,text width=1.5cm] (d) {Dev\\ chave de upload};
        \node[caixapetroleo,text width=1.5cm,right=of d] (g) {Google Play\\ chave do app};
        \node[caixa,right=of g] (u) {Aparelho};
        \draw[seta] (d) -- node[rotulo,above]{AAB} (g);
        \draw[seta] (g) -- node[rotulo,above]{APKs} (u);
      \end{tikzpicture}
      \begin{dica}
        Perdeu a chave de upload? Com o Play App Signing, o Google permite trocá-la. Sem ele, perder a chave impede atualizar o app.
      \end{dica}
    \end{column}
    \begin{column}{0.48\textwidth}
      \begin{conceito}[iOS]
        Todo app precisa ser assinado para rodar num iPhone. Três peças: o \textbf{certificado} (sua identidade), o \textbf{App ID} e o \textbf{perfil de provisionamento}, que liga certificado, app, capacidades e aparelhos.
      \end{conceito}
      \begin{dica}
        O Xcode gerencia tudo isso com \emph{Automatically manage signing}.
      \end{dica}
      \begin{atencao}
        Chaves e certificados são segredos: nunca no Git; guarde cópias num cofre de senhas.
      \end{atencao}
    \end{column}
  \end{columns}
\end{frame}

\begin{frame}[fragile]{Versionamento}
  \begin{columns}[T]
    \begin{column}{0.5\textwidth}
      \begin{definicao}[Versionamento semântico (SemVer)]
        \textbf{MAJOR.MINOR.PATCH}: \textbf{MAJOR} muda com quebras de compatibilidade, \textbf{MINOR} com funcionalidades novas compatíveis e \textbf{PATCH} com correções.
      \end{definicao}
      \centering
      \begin{tikzpicture}[every node/.style={font=\scriptsize}]
        \node[caixapetroleo] (m) {2};
        \node[right=0.5mm of m] (p1) {.};
        \node[caixaambar,right=0.5mm of p1] (n) {4};
        \node[right=0.5mm of n] (p2) {.};
        \node[caixa,right=0.5mm of p2] (c) {1};
        \node[right=0.5mm of c] (p3) {+};
        \node[caixaescura,right=0.5mm of p3] (b) {27};
        \node[rotulo,below=1mm of m,font=\tiny] {major};
        \node[rotulo,below=1mm of n,font=\tiny] {minor};
        \node[rotulo,below=1mm of c,font=\tiny] {patch};
        \node[rotulo,below=1mm of b,font=\tiny,align=center] {número\\do build};
      \end{tikzpicture}
    \end{column}
    \begin{column}{0.48\textwidth}
      \scriptsize
      \begin{tabularx}{\linewidth}{@{}l L L@{}}
        \toprule
        & \cabecalho{Para pessoas} & \cabecalho{Para a loja} \\
        \midrule
        Android & \cod{versionName} & \cod{versionCode} \\
        iOS & \cod{CFBundleShort\-VersionString} & \cod{CFBundleVersion} \\
        Flutter & \cod{version: 2.4.1+27} & (o número após o \cod{+}) \\
        \bottomrule
      \end{tabularx}
      \begin{atencao}[Regra das lojas]
        O número do build \textbf{sempre aumenta}. Não existe ``voltar versão'': para desfazer um problema, publica-se uma versão nova com número maior.
      \end{atencao}
    \end{column}
  \end{columns}
\end{frame}

\begin{frame}{Formatos de pacote e tamanho do app}
  \begin{columns}[T]
    \begin{column}{0.52\textwidth}
      \scriptsize
      \begin{tabularx}{\linewidth}{@{}l L@{}}
        \toprule
        \cabecalho{Formato} & \cabecalho{O que é} \\
        \midrule
        \textbf{APK} & pacote Android instalável, com tudo para todos os aparelhos \\
        \textbf{AAB} & \emph{Android App Bundle}: formato de \textbf{envio} à Play, que gera APKs sob medida para cada aparelho \\
        \textbf{IPA} & pacote iOS enviado à App Store; a Apple também entrega variantes por aparelho (\emph{App Thinning}) \\
        \bottomrule
      \end{tabularx}
      \centering
      \begin{tikzpicture}[node distance=2mm,every node/.style={font=\tiny}]
        \node[caixapetroleo,minimum width=1.6cm] (aab) {AAB};
        \node[caixa,right=8mm of aab,yshift=5mm] (a1) {arm64 + xxhdpi + pt};
        \node[caixa,right=8mm of aab,yshift=-5mm] (a2) {armeabi + hdpi + en};
        \draw[seta] (aab) -- (a1.west); \draw[seta] (aab) -- (a2.west);
      \end{tikzpicture}
    \end{column}
    \begin{column}{0.46\textwidth}
      \begin{conceito}[Por que o tamanho importa]
        Apps grandes demoram a baixar, ocupam espaço e perdem instalações, principalmente em planos de dados limitados.
      \end{conceito}
      \begin{pordentro}[Como os builds encolhem]
        \textbf{\emph{Tree shaking}} e \textbf{\emph{shrinking}} removem código não usado; a \textbf{ofuscação} troca nomes por símbolos curtos; recursos são comprimidos e separados por aparelho.
      \end{pordentro}
      \begin{atencao}
        Código ofuscado gera relatórios de falha ilegíveis. Guarde os \textbf{símbolos de depuração} de cada versão para decodificá-los.
      \end{atencao}
    \end{column}
  \end{columns}
\end{frame}

\begin{frame}{As duas grandes lojas}
  \scriptsize
  \begin{tabularx}{\textwidth}{@{}p{2.6cm} L L@{}}
    \toprule
    & \cabecalho{Google Play} & \cabecalho{App Store} \\
    \midrule
    Conta & taxa única de US\$\,25; pessoal ou organização & Apple Developer Program, US\$\,99 por ano (isenção para instituições de ensino e sem fins lucrativos) \\
    Revisão & automatizada e humana; de horas a dias (contas novas esperam mais) & humana (\emph{App Review}); em geral de um a dois dias \\
    Testes & faixas interna, fechada e aberta & TestFlight interno (até 100) e externo (até 10\,000) \\
    Lançamento gradual & percentual definido por você & \emph{phased release} automático em 7 dias \\
    Requisito anual & \cod{targetSdk} 36 desde 31/08/2026 & Xcode 26 e SDK do iOS 26 desde 28/04/2026 \\
    Comissão & de 15\,\% a 30\,\% sobre vendas digitais, conforme receita, produto e país & idem \\
    \bottomrule
  \end{tabularx}
  \vspace{4pt}
  \begin{dica}
    Regras de comissão e de pagamento externo mudam com leis e decisões judiciais (União Europeia, EUA, Japão, Coreia). Confira antes de definir o modelo de negócio (Módulo 20).
  \end{dica}
\end{frame}

\begin{frame}{Privacidade desde a concepção e LGPD}
  \begin{columns}[T]
    \begin{column}{0.5\textwidth}
      \begin{definicao}[LGPD (Lei 13.709/2018)]
        Regula o tratamento de \textbf{dados pessoais} no Brasil: tudo que identifica ou pode identificar uma pessoa (nome, e-mail, matrícula, localização, identificador do aparelho).
      \end{definicao}
      \begin{conceito}[Princípios que mais afetam um app]
        \textbf{Finalidade} (para quê), \textbf{necessidade} (o mínimo), \textbf{transparência} (explicar), \textbf{segurança} (proteger) e uma \textbf{base legal} para cada uso (consentimento, execução de contrato\ldots).
      \end{conceito}
    \end{column}
    \begin{column}{0.48\textwidth}
      \begin{conceito}[O que as lojas exigem]
        \begin{itemize}
          \item Política de privacidade publicada.
          \item Declaração de dados coletados: \emph{Data safety} (Play) e \emph{App Privacy} (App Store).
          \item Textos de uso para cada permissão.
          \item \textbf{Exclusão de conta} dentro do app, se ele permite criar conta.
        \end{itemize}
      \end{conceito}
      \begin{atencao}
        A declaração deve incluir o que os SDKs de terceiros coletam (analytics, anúncios, Crashlytics). Não é assessoria jurídica: consulte o encarregado de dados (DPO) da instituição.
      \end{atencao}
    \end{column}
  \end{columns}
\end{frame}

\begin{frame}{Gestão de lançamentos}
  \begin{columns}[c]
    \begin{column}{0.52\textwidth}
      \centering
      \begin{tikzpicture}[x=1.0cm,y=0.45cm]
        \draw[cinzamedio,-{Stealth}] (0,0) -- (6.3,0) node[right,font=\tiny]{dias};
        \draw[cinzamedio,-{Stealth}] (0,0) -- (0,5.6) node[above,font=\tiny]{\% dos usuários};
        \draw[azulpetroleo,line width=1.4pt] (0,0.3) -- (1,0.3) -- (1,1) -- (2,1) -- (2,2.5) -- (3.5,2.5) -- (3.5,5) -- (6,5);
        \node[font=\tiny,azulpetroleo] at (0.5,0.9) {1\%};
        \node[font=\tiny,azulpetroleo] at (1.5,1.6) {5\%};
        \node[font=\tiny,azulpetroleo] at (2.75,3.1) {20\%};
        \node[font=\tiny,azulpetroleo] at (4.75,5.6) {100\%};
        \foreach \x in {1,2,3.5} \fill[ambar] (\x,0) circle (2.2pt);
        \node[rotulo,font=\tiny,align=center] at (2.2,-1.6) {\textcolor{ambar}{$\bullet$} portões: falhas e avaliações\\ dentro do limite? então avança};
      \end{tikzpicture}
    \end{column}
    \begin{column}{0.46\textwidth}
      \begin{definicao}[Lançamento gradual]
        Liberar a versão para uma fração crescente dos usuários, avançando só se os indicadores de saúde continuarem bons.
      \end{definicao}
      \begin{conceito}[Ferramentas do processo]
        \begin{itemize}
          \item \textbf{Trem de versões}: lançar em datas fixas (a cada duas semanas, por exemplo).
          \item \textbf{\emph{Feature flags}}: ligar recursos remotamente, sem nova versão.
          \item \textbf{Versão mínima}: obrigar a atualização quando a antiga é insegura.
          \item \textbf{Notas de versão} claras para usuários e revisores.
        \end{itemize}
      \end{conceito}
    \end{column}
  \end{columns}
\end{frame}

\begin{frame}{ASO: o funil da loja}
  \begin{columns}[c]
    \begin{column}{0.46\textwidth}
      \centering
      \begin{tikzpicture}[every node/.style={font=\scriptsize\bfseries,text=white}]
        \fill[azulpetroleo] (-2.6,3.2) -- (2.6,3.2) -- (2.2,2.5) -- (-2.2,2.5) -- cycle;
        \node at (0,2.85) {Impressões na busca};
        \fill[azulprofundo] (-2.15,2.4) -- (2.15,2.4) -- (1.75,1.7) -- (-1.75,1.7) -- cycle;
        \node at (0,2.05) {Visitas à ficha};
        \fill[ambar!90!black] (-1.7,1.6) -- (1.7,1.6) -- (1.3,0.9) -- (-1.3,0.9) -- cycle;
        \node at (0,1.25) {Instalações};
        \fill[ambar] (-1.25,0.8) -- (1.25,0.8) -- (0.85,0.1) -- (-0.85,0.1) -- cycle;
        \node[text=azulpetroleo] at (0,0.45) {Uso ativo};
      \end{tikzpicture}
    \end{column}
    \begin{column}{0.52\textwidth}
      \begin{definicao}[ASO --- \emph{App Store Optimization}]
        Conjunto de práticas para o app ser \textbf{encontrado} na busca da loja e \textbf{escolhido} por quem vê a ficha.
      \end{definicao}
      \scriptsize
      \begin{tabularx}{\linewidth}{@{}l L@{}}
        \toprule
        \cabecalho{Etapa} & \cabecalho{O que influencia} \\
        \midrule
        Ser encontrado & título, subtítulo ou descrição curta, palavras-chave, categoria \\
        Ser escolhido & ícone, capturas de tela, vídeo, nota e avaliações \\
        Ser mantido & qualidade, desempenho, respostas às avaliações \\
        \bottomrule
      \end{tabularx}
      \begin{dica}
        As duas lojas permitem testes A/B da ficha (ícone, capturas, textos).
      \end{dica}
    \end{column}
  \end{columns}
\end{frame}

\begin{frame}{Depois do lançamento: saúde do app}
  \begin{columns}[T]
    \begin{column}{0.5\textwidth}
      \scriptsize
      \begin{tabularx}{\linewidth}{@{}p{2.3cm} L@{}}
        \toprule
        \cabecalho{Indicador} & \cabecalho{O que mede} \\
        \midrule
        Usuários sem falhas & \% de usuários que não tiveram nenhum \emph{crash} no período \\
        Taxa de ANR & \% de sessões em que a tela travou (Android) \\
        Tempo de inicialização & do toque no ícone à primeira tela útil \\
        Retenção D1/D7/D30 & \% de usuários que voltam 1, 7 e 30 dias após instalar \\
        Nota e avaliações & percepção dos usuários; comentários apontam defeitos \\
        \bottomrule
      \end{tabularx}
    \end{column}
    \begin{column}{0.48\textwidth}
      \begin{atencao}[Limites do Android vitals]
        Acima de 1,09\,\% de falhas percebidas ou 0,47\,\% de ANRs percebidos, o app é considerado de ``mau comportamento'' e perde visibilidade na Play.
      \end{atencao}
      \begin{conceito}[Ciclo contínuo]
        Monitorar $\rightarrow$ priorizar $\rightarrow$ corrigir $\rightarrow$ lançar gradualmente $\rightarrow$ monitorar. Falhas e avaliações alimentam o \emph{backlog} da próxima versão.
      \end{conceito}
      \begin{dica}
        Responda às avaliações com educação e objetividade: a resposta é pública e mostra que o app é cuidado.
      \end{dica}
    \end{column}
  \end{columns}
\end{frame}

% =============================================================================
\section{Fundamentos da distribuição de apps}

\begin{frame}{Do código ao aparelho do usuário}
  \centering
  \begin{tikzpicture}[node distance=2.4mm,every node/.style={font=\scriptsize,minimum height=9mm,minimum width=1.62cm}]
    \node[caixa] (c) {Código\\ \tiny versionado};
    \node[caixapetroleo,right=of c] (b) {Build de\\ release};
    \node[caixapetroleo,right=of b] (a) {Assinatura};
    \node[caixaambar,right=of a] (e) {Envio\\ à loja};
    \node[caixaambar,right=of e] (r) {Revisão\\ e políticas};
    \node[caixa,right=of r] (t) {Faixas\\ de teste};
    \node[caixaescura,right=of t] (p) {Produção};
    \foreach \x/\y in {c/b,b/a,a/e,e/r,r/t,t/p} \draw[seta] (\x) -- (\y);
    \draw[setaambar,rounded corners] (p.south) -- ++(0,-6mm) -| node[rotulo,pos=0.25,below]{\tiny monitorar falhas e avaliações $\rightarrow$ nova versão} (c.south);
  \end{tikzpicture}
  \vspace{6pt}
  \begin{columns}[T]
    \begin{column}{0.48\textwidth}
      \scriptsize
      \begin{tabularx}{\linewidth}{@{}l L L@{}}
        \toprule
        & \cabecalho{Debug} & \cabecalho{Release} \\
        \midrule
        Objetivo & desenvolver & distribuir \\
        Código & com verificações e depuração & otimizado, reduzido e ofuscável \\
        Desempenho & mais lento & o real \\
        Assinatura & chave de depuração automática & sua chave de release \\
        \bottomrule
      \end{tabularx}
    \end{column}
    \begin{column}{0.48\textwidth}
      \begin{definicao}[Build de release]
        Versão do app compilada para o usuário final: sem ferramentas de depuração, com código otimizado e assinada com a chave do desenvolvedor.
      \end{definicao}
      \begin{atencao}
        Teste sempre o build de release antes de publicar: ofuscação e otimizações podem expor defeitos que não aparecem em debug.
      \end{atencao}
    \end{column}
  \end{columns}
\end{frame}

\begin{frame}{Assinatura digital: a teoria}
  \begin{columns}[c]
    \begin{column}{0.54\textwidth}
      \centering
      \begin{tikzpicture}[node distance=4mm,every node/.style={font=\scriptsize}]
        \node[caixa,minimum width=1.6cm] (app) {App};
        \node[caixa,right=of app,minimum width=1.6cm] (h) {Hash\\ \tiny resumo único};
        \node[caixaescura,right=of h,minimum width=2cm] (s) {Assinatura};
        \node[caixaambar,above=5mm of s,minimum width=2cm] (pk) {Chave \textbf{privada}\\ \tiny só o dev tem};
        \draw[seta] (app) -- (h); \draw[seta] (h) -- (s); \draw[seta] (pk) -- (s);
        \node[caixa,below=9mm of app,minimum width=1.6cm] (app2) {App\\ \tiny recebido};
        \node[caixa,right=of app2,minimum width=1.6cm] (h2) {Hash\\ \tiny recalculado};
        \node[caixapetroleo,right=of h2,minimum width=2cm] (v) {Confere?};
        \node[caixa,below=5mm of v,minimum width=2cm] (pub) {Chave \textbf{pública}\\ \tiny no certificado};
        \draw[seta] (app2) -- (h2); \draw[seta] (h2) -- (v); \draw[seta] (pub) -- (v);
        \draw[setaambar] (s.south) -- node[rotulo,right]{\tiny vai junto} (v.north);
        \node[rotulo,left=1mm of app,font=\tiny,rotate=90,anchor=south] {quem publica};
        \node[rotulo,left=1mm of app2,font=\tiny,rotate=90,anchor=south] {sistema};
      \end{tikzpicture}
    \end{column}
    \begin{column}{0.44\textwidth}
      \begin{definicao}[Assinatura digital]
        Resumo (\emph{hash}) do pacote cifrado com a \textbf{chave privada} do desenvolvedor. Qualquer um verifica com a \textbf{chave pública}, mas só quem tem a privada consegue assinar.
      \end{definicao}
      \begin{conceito}[O que ela garante]
        \textbf{Integridade}: o pacote não foi alterado.\\
        \textbf{Autenticidade}: a atualização vem do mesmo autor do app instalado.
      \end{conceito}
      \begin{analogia}
        O lacre de cera com o sinete do remetente: qualquer um reconhece o selo, ninguém o falsifica sem o anel.
      \end{analogia}
    \end{column}
  \end{columns}
\end{frame}

\begin{frame}{Assinatura em cada plataforma}
  \begin{columns}[T]
    \begin{column}{0.5\textwidth}
      \begin{conceito}[Android]
        O pacote é assinado com uma chave guardada num \emph{keystore}. Com o \textbf{Play App Signing}, você assina com a \textbf{chave de upload}; o Google confere e reassina com a \textbf{chave do app}, que fica guardada com ele.
      \end{conceito}
      \centering
      \begin{tikzpicture}[node distance=9mm,every node/.style={font=\tiny,align=center}]
        \node[caixa,text width=1.5cm] (d) {Dev\\ chave de upload};
        \node[caixapetroleo,text width=1.5cm,right=of d] (g) {Google Play\\ chave do app};
        \node[caixa,right=of g] (u) {Aparelho};
        \draw[seta] (d) -- node[rotulo,above]{AAB} (g);
        \draw[seta] (g) -- node[rotulo,above]{APKs} (u);
      \end{tikzpicture}
      \begin{dica}
        Perdeu a chave de upload? Com o Play App Signing, o Google permite trocá-la. Sem ele, perder a chave impede atualizar o app.
      \end{dica}
    \end{column}
    \begin{column}{0.48\textwidth}
      \begin{conceito}[iOS]
        Todo app precisa ser assinado para rodar num iPhone. Três peças: o \textbf{certificado} (sua identidade), o \textbf{App ID} e o \textbf{perfil de provisionamento}, que liga certificado, app, capacidades e aparelhos.
      \end{conceito}
      \begin{dica}
        O Xcode gerencia tudo isso com \emph{Automatically manage signing}.
      \end{dica}
      \begin{atencao}
        Chaves e certificados são segredos: nunca no Git; guarde cópias num cofre de senhas.
      \end{atencao}
    \end{column}
  \end{columns}
\end{frame}

\begin{frame}[fragile]{Versionamento}
  \begin{columns}[T]
    \begin{column}{0.5\textwidth}
      \begin{definicao}[Versionamento semântico (SemVer)]
        \textbf{MAJOR.MINOR.PATCH}: \textbf{MAJOR} muda com quebras de compatibilidade, \textbf{MINOR} com funcionalidades novas compatíveis e \textbf{PATCH} com correções.
      \end{definicao}
      \centering
      \begin{tikzpicture}[every node/.style={font=\scriptsize}]
        \node[caixapetroleo] (m) {2};
        \node[right=0.5mm of m] (p1) {.};
        \node[caixaambar,right=0.5mm of p1] (n) {4};
        \node[right=0.5mm of n] (p2) {.};
        \node[caixa,right=0.5mm of p2] (c) {1};
        \node[right=0.5mm of c] (p3) {+};
        \node[caixaescura,right=0.5mm of p3] (b) {27};
        \node[rotulo,below=1mm of m,font=\tiny] {major};
        \node[rotulo,below=1mm of n,font=\tiny] {minor};
        \node[rotulo,below=1mm of c,font=\tiny] {patch};
        \node[rotulo,below=1mm of b,font=\tiny,align=center] {número\\do build};
      \end{tikzpicture}
    \end{column}
    \begin{column}{0.48\textwidth}
      \scriptsize
      \begin{tabularx}{\linewidth}{@{}l L L@{}}
        \toprule
        & \cabecalho{Para pessoas} & \cabecalho{Para a loja} \\
        \midrule
        Android & \cod{versionName} & \cod{versionCode} \\
        iOS & \cod{CFBundleShort\-VersionString} & \cod{CFBundleVersion} \\
        Flutter & \cod{version: 2.4.1+27} & (o número após o \cod{+}) \\
        \bottomrule
      \end{tabularx}
      \begin{atencao}[Regra das lojas]
        O número do build \textbf{sempre aumenta}. Não existe ``voltar versão'': para desfazer um problema, publica-se uma versão nova com número maior.
      \end{atencao}
    \end{column}
  \end{columns}
\end{frame}

\begin{frame}{Formatos de pacote e tamanho do app}
  \begin{columns}[T]
    \begin{column}{0.52\textwidth}
      \scriptsize
      \begin{tabularx}{\linewidth}{@{}l L@{}}
        \toprule
        \cabecalho{Formato} & \cabecalho{O que é} \\
        \midrule
        \textbf{APK} & pacote Android instalável, com tudo para todos os aparelhos \\
        \textbf{AAB} & \emph{Android App Bundle}: formato de \textbf{envio} à Play, que gera APKs sob medida para cada aparelho \\
        \textbf{IPA} & pacote iOS enviado à App Store; a Apple também entrega variantes por aparelho (\emph{App Thinning}) \\
        \bottomrule
      \end{tabularx}
      \centering
      \begin{tikzpicture}[node distance=2mm,every node/.style={font=\tiny}]
        \node[caixapetroleo,minimum width=1.6cm] (aab) {AAB};
        \node[caixa,right=8mm of aab,yshift=5mm] (a1) {arm64 + xxhdpi + pt};
        \node[caixa,right=8mm of aab,yshift=-5mm] (a2) {armeabi + hdpi + en};
        \draw[seta] (aab) -- (a1.west); \draw[seta] (aab) -- (a2.west);
      \end{tikzpicture}
    \end{column}
    \begin{column}{0.46\textwidth}
      \begin{conceito}[Por que o tamanho importa]
        Apps grandes demoram a baixar, ocupam espaço e perdem instalações, principalmente em planos de dados limitados.
      \end{conceito}
      \begin{pordentro}[Como os builds encolhem]
        \textbf{\emph{Tree shaking}} e \textbf{\emph{shrinking}} removem código não usado; a \textbf{ofuscação} troca nomes por símbolos curtos; recursos são comprimidos e separados por aparelho.
      \end{pordentro}
      \begin{atencao}
        Código ofuscado gera relatórios de falha ilegíveis. Guarde os \textbf{símbolos de depuração} de cada versão para decodificá-los.
      \end{atencao}
    \end{column}
  \end{columns}
\end{frame}

\begin{frame}{As duas grandes lojas}
  \scriptsize
  \begin{tabularx}{\textwidth}{@{}p{2.6cm} L L@{}}
    \toprule
    & \cabecalho{Google Play} & \cabecalho{App Store} \\
    \midrule
    Conta & taxa única de US\$\,25; pessoal ou organização & Apple Developer Program, US\$\,99 por ano (isenção para instituições de ensino e sem fins lucrativos) \\
    Revisão & automatizada e humana; de horas a dias (contas novas esperam mais) & humana (\emph{App Review}); em geral de um a dois dias \\
    Testes & faixas interna, fechada e aberta & TestFlight interno (até 100) e externo (até 10\,000) \\
    Lançamento gradual & percentual definido por você & \emph{phased release} automático em 7 dias \\
    Requisito anual & \cod{targetSdk} 36 desde 31/08/2026 & Xcode 26 e SDK do iOS 26 desde 28/04/2026 \\
    Comissão & de 15\,\% a 30\,\% sobre vendas digitais, conforme receita, produto e país & idem \\
    \bottomrule
  \end{tabularx}
  \vspace{4pt}
  \begin{dica}
    Regras de comissão e de pagamento externo mudam com leis e decisões judiciais (União Europeia, EUA, Japão, Coreia). Confira antes de definir o modelo de negócio (Módulo 20).
  \end{dica}
\end{frame}

\begin{frame}{Privacidade desde a concepção e LGPD}
  \begin{columns}[T]
    \begin{column}{0.5\textwidth}
      \begin{definicao}[LGPD (Lei 13.709/2018)]
        Regula o tratamento de \textbf{dados pessoais} no Brasil: tudo que identifica ou pode identificar uma pessoa (nome, e-mail, matrícula, localização, identificador do aparelho).
      \end{definicao}
      \begin{conceito}[Princípios que mais afetam um app]
        \textbf{Finalidade} (para quê), \textbf{necessidade} (o mínimo), \textbf{transparência} (explicar), \textbf{segurança} (proteger) e uma \textbf{base legal} para cada uso (consentimento, execução de contrato\ldots).
      \end{conceito}
    \end{column}
    \begin{column}{0.48\textwidth}
      \begin{conceito}[O que as lojas exigem]
        \begin{itemize}
          \item Política de privacidade publicada.
          \item Declaração de dados coletados: \emph{Data safety} (Play) e \emph{App Privacy} (App Store).
          \item Textos de uso para cada permissão.
          \item \textbf{Exclusão de conta} dentro do app, se ele permite criar conta.
        \end{itemize}
      \end{conceito}
      \begin{atencao}
        A declaração deve incluir o que os SDKs de terceiros coletam (analytics, anúncios, Crashlytics). Não é assessoria jurídica: consulte o encarregado de dados (DPO) da instituição.
      \end{atencao}
    \end{column}
  \end{columns}
\end{frame}

\begin{frame}{Gestão de lançamentos}
  \begin{columns}[c]
    \begin{column}{0.52\textwidth}
      \centering
      \begin{tikzpicture}[x=1.0cm,y=0.45cm]
        \draw[cinzamedio,-{Stealth}] (0,0) -- (6.3,0) node[right,font=\tiny]{dias};
        \draw[cinzamedio,-{Stealth}] (0,0) -- (0,5.6) node[above,font=\tiny]{\% dos usuários};
        \draw[azulpetroleo,line width=1.4pt] (0,0.3) -- (1,0.3) -- (1,1) -- (2,1) -- (2,2.5) -- (3.5,2.5) -- (3.5,5) -- (6,5);
        \node[font=\tiny,azulpetroleo] at (0.5,0.9) {1\%};
        \node[font=\tiny,azulpetroleo] at (1.5,1.6) {5\%};
        \node[font=\tiny,azulpetroleo] at (2.75,3.1) {20\%};
        \node[font=\tiny,azulpetroleo] at (4.75,5.6) {100\%};
        \foreach \x in {1,2,3.5} \fill[ambar] (\x,0) circle (2.2pt);
        \node[rotulo,font=\tiny,align=center] at (2.2,-1.6) {\textcolor{ambar}{$\bullet$} portões: falhas e avaliações\\ dentro do limite? então avança};
      \end{tikzpicture}
    \end{column}
    \begin{column}{0.46\textwidth}
      \begin{definicao}[Lançamento gradual]
        Liberar a versão para uma fração crescente dos usuários, avançando só se os indicadores de saúde continuarem bons.
      \end{definicao}
      \begin{conceito}[Ferramentas do processo]
        \begin{itemize}
          \item \textbf{Trem de versões}: lançar em datas fixas (a cada duas semanas, por exemplo).
          \item \textbf{\emph{Feature flags}}: ligar recursos remotamente, sem nova versão.
          \item \textbf{Versão mínima}: obrigar a atualização quando a antiga é insegura.
          \item \textbf{Notas de versão} claras para usuários e revisores.
        \end{itemize}
      \end{conceito}
    \end{column}
  \end{columns}
\end{frame}

\begin{frame}{ASO: o funil da loja}
  \begin{columns}[c]
    \begin{column}{0.46\textwidth}
      \centering
      \begin{tikzpicture}[every node/.style={font=\scriptsize\bfseries,text=white}]
        \fill[azulpetroleo] (-2.6,3.2) -- (2.6,3.2) -- (2.2,2.5) -- (-2.2,2.5) -- cycle;
        \node at (0,2.85) {Impressões na busca};
        \fill[azulprofundo] (-2.15,2.4) -- (2.15,2.4) -- (1.75,1.7) -- (-1.75,1.7) -- cycle;
        \node at (0,2.05) {Visitas à ficha};
        \fill[ambar!90!black] (-1.7,1.6) -- (1.7,1.6) -- (1.3,0.9) -- (-1.3,0.9) -- cycle;
        \node at (0,1.25) {Instalações};
        \fill[ambar] (-1.25,0.8) -- (1.25,0.8) -- (0.85,0.1) -- (-0.85,0.1) -- cycle;
        \node[text=azulpetroleo] at (0,0.45) {Uso ativo};
      \end{tikzpicture}
    \end{column}
    \begin{column}{0.52\textwidth}
      \begin{definicao}[ASO --- \emph{App Store Optimization}]
        Conjunto de práticas para o app ser \textbf{encontrado} na busca da loja e \textbf{escolhido} por quem vê a ficha.
      \end{definicao}
      \scriptsize
      \begin{tabularx}{\linewidth}{@{}l L@{}}
        \toprule
        \cabecalho{Etapa} & \cabecalho{O que influencia} \\
        \midrule
        Ser encontrado & título, subtítulo ou descrição curta, palavras-chave, categoria \\
        Ser escolhido & ícone, capturas de tela, vídeo, nota e avaliações \\
        Ser mantido & qualidade, desempenho, respostas às avaliações \\
        \bottomrule
      \end{tabularx}
      \begin{dica}
        As duas lojas permitem testes A/B da ficha (ícone, capturas, textos).
      \end{dica}
    \end{column}
  \end{columns}
\end{frame}

\begin{frame}{Depois do lançamento: saúde do app}
  \begin{columns}[T]
    \begin{column}{0.5\textwidth}
      \scriptsize
      \begin{tabularx}{\linewidth}{@{}p{2.3cm} L@{}}
        \toprule
        \cabecalho{Indicador} & \cabecalho{O que mede} \\
        \midrule
        Usuários sem falhas & \% de usuários que não tiveram nenhum \emph{crash} no período \\
        Taxa de ANR & \% de sessões em que a tela travou (Android) \\
        Tempo de inicialização & do toque no ícone à primeira tela útil \\
        Retenção D1/D7/D30 & \% de usuários que voltam 1, 7 e 30 dias após instalar \\
        Nota e avaliações & percepção dos usuários; comentários apontam defeitos \\
        \bottomrule
      \end{tabularx}
    \end{column}
    \begin{column}{0.48\textwidth}
      \begin{atencao}[Limites do Android vitals]
        Acima de 1,09\,\% de falhas percebidas ou 0,47\,\% de ANRs percebidos, o app é considerado de ``mau comportamento'' e perde visibilidade na Play.
      \end{atencao}
      \begin{conceito}[Ciclo contínuo]
        Monitorar $\rightarrow$ priorizar $\rightarrow$ corrigir $\rightarrow$ lançar gradualmente $\rightarrow$ monitorar. Falhas e avaliações alimentam o \emph{backlog} da próxima versão.
      \end{conceito}
      \begin{dica}
        Responda às avaliações com educação e objetividade: a resposta é pública e mostra que o app é cuidado.
      \end{dica}
    \end{column}
  \end{columns}
\end{frame}

% =============================================================================
\section{Fundamentos da distribuição de apps}

\begin{frame}{Do código ao aparelho do usuário}
  \centering
  \begin{tikzpicture}[node distance=2.4mm,every node/.style={font=\scriptsize,minimum height=9mm,minimum width=1.62cm}]
    \node[caixa] (c) {Código\\ \tiny versionado};
    \node[caixapetroleo,right=of c] (b) {Build de\\ release};
    \node[caixapetroleo,right=of b] (a) {Assinatura};
    \node[caixaambar,right=of a] (e) {Envio\\ à loja};
    \node[caixaambar,right=of e] (r) {Revisão\\ e políticas};
    \node[caixa,right=of r] (t) {Faixas\\ de teste};
    \node[caixaescura,right=of t] (p) {Produção};
    \foreach \x/\y in {c/b,b/a,a/e,e/r,r/t,t/p} \draw[seta] (\x) -- (\y);
    \draw[setaambar,rounded corners] (p.south) -- ++(0,-6mm) -| node[rotulo,pos=0.25,below]{\tiny monitorar falhas e avaliações $\rightarrow$ nova versão} (c.south);
  \end{tikzpicture}
  \vspace{6pt}
  \begin{columns}[T]
    \begin{column}{0.48\textwidth}
      \scriptsize
      \begin{tabularx}{\linewidth}{@{}l L L@{}}
        \toprule
        & \cabecalho{Debug} & \cabecalho{Release} \\
        \midrule
        Objetivo & desenvolver & distribuir \\
        Código & com verificações e depuração & otimizado, reduzido e ofuscável \\
        Desempenho & mais lento & o real \\
        Assinatura & chave de depuração automática & sua chave de release \\
        \bottomrule
      \end{tabularx}
    \end{column}
    \begin{column}{0.48\textwidth}
      \begin{definicao}[Build de release]
        Versão do app compilada para o usuário final: sem ferramentas de depuração, com código otimizado e assinada com a chave do desenvolvedor.
      \end{definicao}
      \begin{atencao}
        Teste sempre o build de release antes de publicar: ofuscação e otimizações podem expor defeitos que não aparecem em debug.
      \end{atencao}
    \end{column}
  \end{columns}
\end{frame}

\begin{frame}{Assinatura digital: a teoria}
  \begin{columns}[c]
    \begin{column}{0.54\textwidth}
      \centering
      \begin{tikzpicture}[node distance=4mm,every node/.style={font=\scriptsize}]
        \node[caixa,minimum width=1.6cm] (app) {App};
        \node[caixa,right=of app,minimum width=1.6cm] (h) {Hash\\ \tiny resumo único};
        \node[caixaescura,right=of h,minimum width=2cm] (s) {Assinatura};
        \node[caixaambar,above=5mm of s,minimum width=2cm] (pk) {Chave \textbf{privada}\\ \tiny só o dev tem};
        \draw[seta] (app) -- (h); \draw[seta] (h) -- (s); \draw[seta] (pk) -- (s);
        \node[caixa,below=9mm of app,minimum width=1.6cm] (app2) {App\\ \tiny recebido};
        \node[caixa,right=of app2,minimum width=1.6cm] (h2) {Hash\\ \tiny recalculado};
        \node[caixapetroleo,right=of h2,minimum width=2cm] (v) {Confere?};
        \node[caixa,below=5mm of v,minimum width=2cm] (pub) {Chave \textbf{pública}\\ \tiny no certificado};
        \draw[seta] (app2) -- (h2); \draw[seta] (h2) -- (v); \draw[seta] (pub) -- (v);
        \draw[setaambar] (s.south) -- node[rotulo,right]{\tiny vai junto} (v.north);
        \node[rotulo,left=1mm of app,font=\tiny,rotate=90,anchor=south] {quem publica};
        \node[rotulo,left=1mm of app2,font=\tiny,rotate=90,anchor=south] {sistema};
      \end{tikzpicture}
    \end{column}
    \begin{column}{0.44\textwidth}
      \begin{definicao}[Assinatura digital]
        Resumo (\emph{hash}) do pacote cifrado com a \textbf{chave privada} do desenvolvedor. Qualquer um verifica com a \textbf{chave pública}, mas só quem tem a privada consegue assinar.
      \end{definicao}
      \begin{conceito}[O que ela garante]
        \textbf{Integridade}: o pacote não foi alterado.\\
        \textbf{Autenticidade}: a atualização vem do mesmo autor do app instalado.
      \end{conceito}
      \begin{analogia}
        O lacre de cera com o sinete do remetente: qualquer um reconhece o selo, ninguém o falsifica sem o anel.
      \end{analogia}
    \end{column}
  \end{columns}
\end{frame}

\begin{frame}{Assinatura em cada plataforma}
  \begin{columns}[T]
    \begin{column}{0.5\textwidth}
      \begin{conceito}[Android]
        O pacote é assinado com uma chave guardada num \emph{keystore}. Com o \textbf{Play App Signing}, você assina com a \textbf{chave de upload}; o Google confere e reassina com a \textbf{chave do app}, que fica guardada com ele.
      \end{conceito}
      \centering
      \begin{tikzpicture}[node distance=9mm,every node/.style={font=\tiny,align=center}]
        \node[caixa,text width=1.5cm] (d) {Dev\\ chave de upload};
        \node[caixapetroleo,text width=1.5cm,right=of d] (g) {Google Play\\ chave do app};
        \node[caixa,right=of g] (u) {Aparelho};
        \draw[seta] (d) -- node[rotulo,above]{AAB} (g);
        \draw[seta] (g) -- node[rotulo,above]{APKs} (u);
      \end{tikzpicture}
      \begin{dica}
        Perdeu a chave de upload? Com o Play App Signing, o Google permite trocá-la. Sem ele, perder a chave impede atualizar o app.
      \end{dica}
    \end{column}
    \begin{column}{0.48\textwidth}
      \begin{conceito}[iOS]
        Todo app precisa ser assinado para rodar num iPhone. Três peças: o \textbf{certificado} (sua identidade), o \textbf{App ID} e o \textbf{perfil de provisionamento}, que liga certificado, app, capacidades e aparelhos.
      \end{conceito}
      \begin{dica}
        O Xcode gerencia tudo isso com \emph{Automatically manage signing}.
      \end{dica}
      \begin{atencao}
        Chaves e certificados são segredos: nunca no Git; guarde cópias num cofre de senhas.
      \end{atencao}
    \end{column}
  \end{columns}
\end{frame}

\begin{frame}[fragile]{Versionamento}
  \begin{columns}[T]
    \begin{column}{0.5\textwidth}
      \begin{definicao}[Versionamento semântico (SemVer)]
        \textbf{MAJOR.MINOR.PATCH}: \textbf{MAJOR} muda com quebras de compatibilidade, \textbf{MINOR} com funcionalidades novas compatíveis e \textbf{PATCH} com correções.
      \end{definicao}
      \centering
      \begin{tikzpicture}[every node/.style={font=\scriptsize}]
        \node[caixapetroleo] (m) {2};
        \node[right=0.5mm of m] (p1) {.};
        \node[caixaambar,right=0.5mm of p1] (n) {4};
        \node[right=0.5mm of n] (p2) {.};
        \node[caixa,right=0.5mm of p2] (c) {1};
        \node[right=0.5mm of c] (p3) {+};
        \node[caixaescura,right=0.5mm of p3] (b) {27};
        \node[rotulo,below=1mm of m,font=\tiny] {major};
        \node[rotulo,below=1mm of n,font=\tiny] {minor};
        \node[rotulo,below=1mm of c,font=\tiny] {patch};
        \node[rotulo,below=1mm of b,font=\tiny,align=center] {número\\do build};
      \end{tikzpicture}
    \end{column}
    \begin{column}{0.48\textwidth}
      \scriptsize
      \begin{tabularx}{\linewidth}{@{}l L L@{}}
        \toprule
        & \cabecalho{Para pessoas} & \cabecalho{Para a loja} \\
        \midrule
        Android & \cod{versionName} & \cod{versionCode} \\
        iOS & \cod{CFBundleShort\-VersionString} & \cod{CFBundleVersion} \\
        Flutter & \cod{version: 2.4.1+27} & (o número após o \cod{+}) \\
        \bottomrule
      \end{tabularx}
      \begin{atencao}[Regra das lojas]
        O número do build \textbf{sempre aumenta}. Não existe ``voltar versão'': para desfazer um problema, publica-se uma versão nova com número maior.
      \end{atencao}
    \end{column}
  \end{columns}
\end{frame}

\begin{frame}{Formatos de pacote e tamanho do app}
  \begin{columns}[T]
    \begin{column}{0.52\textwidth}
      \scriptsize
      \begin{tabularx}{\linewidth}{@{}l L@{}}
        \toprule
        \cabecalho{Formato} & \cabecalho{O que é} \\
        \midrule
        \textbf{APK} & pacote Android instalável, com tudo para todos os aparelhos \\
        \textbf{AAB} & \emph{Android App Bundle}: formato de \textbf{envio} à Play, que gera APKs sob medida para cada aparelho \\
        \textbf{IPA} & pacote iOS enviado à App Store; a Apple também entrega variantes por aparelho (\emph{App Thinning}) \\
        \bottomrule
      \end{tabularx}
      \centering
      \begin{tikzpicture}[node distance=2mm,every node/.style={font=\tiny}]
        \node[caixapetroleo,minimum width=1.6cm] (aab) {AAB};
        \node[caixa,right=8mm of aab,yshift=5mm] (a1) {arm64 + xxhdpi + pt};
        \node[caixa,right=8mm of aab,yshift=-5mm] (a2) {armeabi + hdpi + en};
        \draw[seta] (aab) -- (a1.west); \draw[seta] (aab) -- (a2.west);
      \end{tikzpicture}
    \end{column}
    \begin{column}{0.46\textwidth}
      \begin{conceito}[Por que o tamanho importa]
        Apps grandes demoram a baixar, ocupam espaço e perdem instalações, principalmente em planos de dados limitados.
      \end{conceito}
      \begin{pordentro}[Como os builds encolhem]
        \textbf{\emph{Tree shaking}} e \textbf{\emph{shrinking}} removem código não usado; a \textbf{ofuscação} troca nomes por símbolos curtos; recursos são comprimidos e separados por aparelho.
      \end{pordentro}
      \begin{atencao}
        Código ofuscado gera relatórios de falha ilegíveis. Guarde os \textbf{símbolos de depuração} de cada versão para decodificá-los.
      \end{atencao}
    \end{column}
  \end{columns}
\end{frame}

\begin{frame}{As duas grandes lojas}
  \scriptsize
  \begin{tabularx}{\textwidth}{@{}p{2.6cm} L L@{}}
    \toprule
    & \cabecalho{Google Play} & \cabecalho{App Store} \\
    \midrule
    Conta & taxa única de US\$\,25; pessoal ou organização & Apple Developer Program, US\$\,99 por ano (isenção para instituições de ensino e sem fins lucrativos) \\
    Revisão & automatizada e humana; de horas a dias (contas novas esperam mais) & humana (\emph{App Review}); em geral de um a dois dias \\
    Testes & faixas interna, fechada e aberta & TestFlight interno (até 100) e externo (até 10\,000) \\
    Lançamento gradual & percentual definido por você & \emph{phased release} automático em 7 dias \\
    Requisito anual & \cod{targetSdk} 36 desde 31/08/2026 & Xcode 26 e SDK do iOS 26 desde 28/04/2026 \\
    Comissão & de 15\,\% a 30\,\% sobre vendas digitais, conforme receita, produto e país & idem \\
    \bottomrule
  \end{tabularx}
  \vspace{4pt}
  \begin{dica}
    Regras de comissão e de pagamento externo mudam com leis e decisões judiciais (União Europeia, EUA, Japão, Coreia). Confira antes de definir o modelo de negócio (Módulo 20).
  \end{dica}
\end{frame}

\begin{frame}{Privacidade desde a concepção e LGPD}
  \begin{columns}[T]
    \begin{column}{0.5\textwidth}
      \begin{definicao}[LGPD (Lei 13.709/2018)]
        Regula o tratamento de \textbf{dados pessoais} no Brasil: tudo que identifica ou pode identificar uma pessoa (nome, e-mail, matrícula, localização, identificador do aparelho).
      \end{definicao}
      \begin{conceito}[Princípios que mais afetam um app]
        \textbf{Finalidade} (para quê), \textbf{necessidade} (o mínimo), \textbf{transparência} (explicar), \textbf{segurança} (proteger) e uma \textbf{base legal} para cada uso (consentimento, execução de contrato\ldots).
      \end{conceito}
    \end{column}
    \begin{column}{0.48\textwidth}
      \begin{conceito}[O que as lojas exigem]
        \begin{itemize}
          \item Política de privacidade publicada.
          \item Declaração de dados coletados: \emph{Data safety} (Play) e \emph{App Privacy} (App Store).
          \item Textos de uso para cada permissão.
          \item \textbf{Exclusão de conta} dentro do app, se ele permite criar conta.
        \end{itemize}
      \end{conceito}
      \begin{atencao}
        A declaração deve incluir o que os SDKs de terceiros coletam (analytics, anúncios, Crashlytics). Não é assessoria jurídica: consulte o encarregado de dados (DPO) da instituição.
      \end{atencao}
    \end{column}
  \end{columns}
\end{frame}

\begin{frame}{Gestão de lançamentos}
  \begin{columns}[c]
    \begin{column}{0.52\textwidth}
      \centering
      \begin{tikzpicture}[x=1.0cm,y=0.45cm]
        \draw[cinzamedio,-{Stealth}] (0,0) -- (6.3,0) node[right,font=\tiny]{dias};
        \draw[cinzamedio,-{Stealth}] (0,0) -- (0,5.6) node[above,font=\tiny]{\% dos usuários};
        \draw[azulpetroleo,line width=1.4pt] (0,0.3) -- (1,0.3) -- (1,1) -- (2,1) -- (2,2.5) -- (3.5,2.5) -- (3.5,5) -- (6,5);
        \node[font=\tiny,azulpetroleo] at (0.5,0.9) {1\%};
        \node[font=\tiny,azulpetroleo] at (1.5,1.6) {5\%};
        \node[font=\tiny,azulpetroleo] at (2.75,3.1) {20\%};
        \node[font=\tiny,azulpetroleo] at (4.75,5.6) {100\%};
        \foreach \x in {1,2,3.5} \fill[ambar] (\x,0) circle (2.2pt);
        \node[rotulo,font=\tiny,align=center] at (2.2,-1.6) {\textcolor{ambar}{$\bullet$} portões: falhas e avaliações\\ dentro do limite? então avança};
      \end{tikzpicture}
    \end{column}
    \begin{column}{0.46\textwidth}
      \begin{definicao}[Lançamento gradual]
        Liberar a versão para uma fração crescente dos usuários, avançando só se os indicadores de saúde continuarem bons.
      \end{definicao}
      \begin{conceito}[Ferramentas do processo]
        \begin{itemize}
          \item \textbf{Trem de versões}: lançar em datas fixas (a cada duas semanas, por exemplo).
          \item \textbf{\emph{Feature flags}}: ligar recursos remotamente, sem nova versão.
          \item \textbf{Versão mínima}: obrigar a atualização quando a antiga é insegura.
          \item \textbf{Notas de versão} claras para usuários e revisores.
        \end{itemize}
      \end{conceito}
    \end{column}
  \end{columns}
\end{frame}

\begin{frame}{ASO: o funil da loja}
  \begin{columns}[c]
    \begin{column}{0.46\textwidth}
      \centering
      \begin{tikzpicture}[every node/.style={font=\scriptsize\bfseries,text=white}]
        \fill[azulpetroleo] (-2.6,3.2) -- (2.6,3.2) -- (2.2,2.5) -- (-2.2,2.5) -- cycle;
        \node at (0,2.85) {Impressões na busca};
        \fill[azulprofundo] (-2.15,2.4) -- (2.15,2.4) -- (1.75,1.7) -- (-1.75,1.7) -- cycle;
        \node at (0,2.05) {Visitas à ficha};
        \fill[ambar!90!black] (-1.7,1.6) -- (1.7,1.6) -- (1.3,0.9) -- (-1.3,0.9) -- cycle;
        \node at (0,1.25) {Instalações};
        \fill[ambar] (-1.25,0.8) -- (1.25,0.8) -- (0.85,0.1) -- (-0.85,0.1) -- cycle;
        \node[text=azulpetroleo] at (0,0.45) {Uso ativo};
      \end{tikzpicture}
    \end{column}
    \begin{column}{0.52\textwidth}
      \begin{definicao}[ASO --- \emph{App Store Optimization}]
        Conjunto de práticas para o app ser \textbf{encontrado} na busca da loja e \textbf{escolhido} por quem vê a ficha.
      \end{definicao}
      \scriptsize
      \begin{tabularx}{\linewidth}{@{}l L@{}}
        \toprule
        \cabecalho{Etapa} & \cabecalho{O que influencia} \\
        \midrule
        Ser encontrado & título, subtítulo ou descrição curta, palavras-chave, categoria \\
        Ser escolhido & ícone, capturas de tela, vídeo, nota e avaliações \\
        Ser mantido & qualidade, desempenho, respostas às avaliações \\
        \bottomrule
      \end{tabularx}
      \begin{dica}
        As duas lojas permitem testes A/B da ficha (ícone, capturas, textos).
      \end{dica}
    \end{column}
  \end{columns}
\end{frame}

\begin{frame}{Depois do lançamento: saúde do app}
  \begin{columns}[T]
    \begin{column}{0.5\textwidth}
      \scriptsize
      \begin{tabularx}{\linewidth}{@{}p{2.3cm} L@{}}
        \toprule
        \cabecalho{Indicador} & \cabecalho{O que mede} \\
        \midrule
        Usuários sem falhas & \% de usuários que não tiveram nenhum \emph{crash} no período \\
        Taxa de ANR & \% de sessões em que a tela travou (Android) \\
        Tempo de inicialização & do toque no ícone à primeira tela útil \\
        Retenção D1/D7/D30 & \% de usuários que voltam 1, 7 e 30 dias após instalar \\
        Nota e avaliações & percepção dos usuários; comentários apontam defeitos \\
        \bottomrule
      \end{tabularx}
    \end{column}
    \begin{column}{0.48\textwidth}
      \begin{atencao}[Limites do Android vitals]
        Acima de 1,09\,\% de falhas percebidas ou 0,47\,\% de ANRs percebidos, o app é considerado de ``mau comportamento'' e perde visibilidade na Play.
      \end{atencao}
      \begin{conceito}[Ciclo contínuo]
        Monitorar $\rightarrow$ priorizar $\rightarrow$ corrigir $\rightarrow$ lançar gradualmente $\rightarrow$ monitorar. Falhas e avaliações alimentam o \emph{backlog} da próxima versão.
      \end{conceito}
      \begin{dica}
        Responda às avaliações com educação e objetividade: a resposta é pública e mostra que o app é cuidado.
      \end{dica}
    \end{column}
  \end{columns}
\end{frame}


% =============================================================================
\section{Preparando o release}

\begin{frame}[fragile]{Build de release}
  \begin{columns}[T]
    \begin{column}{0.54\textwidth}
\begin{codigo}[title=app/build.gradle.kts]{Kotlin}
android {
  defaultConfig {
    versionCode = 7          // inteiro, sempre crescente
    versionName = "1.2.0"    // exibido ao usuário
  }
  buildTypes {
    release {
      isMinifyEnabled = true      // R8: encolhe e ofusca
      isShrinkResources = true    // remove recursos não usados
      proguardFiles(getDefaultProguardFile(
          "proguard-android-optimize.txt"),
        "proguard-rules.pro")
      signingConfig = signingConfigs.getByName("release")
    }
  }
}
\end{codigo}
    \end{column}
    \begin{column}{0.44\textwidth}
      \begin{conceito}[Versionamento]
        \cod{versionCode} cresce a cada envio. \cod{versionName} segue o versionamento semântico (maior.menor.correção).
      \end{conceito}
      \begin{atencao}[R8]
        Renomeia classes: bibliotecas com reflexão podem exigir regras \cod{-keep}. \textbf{Teste o build release.}
      \end{atencao}
      \begin{dica}
        Guarde o \arq{mapping.txt} de cada versão: ele desofusca as pilhas de erro no Crashlytics e na Play Console.
      \end{dica}
    \end{column}
  \end{columns}
\end{frame}

\begin{frame}{Checklist antes de publicar}
  \begin{columns}[T]
    \begin{column}{0.48\textwidth}
      \begin{conceito}[Produto]
        \begin{itemize}
          \item Ícone adaptativo (\emph{Image Asset Studio}) e nome definitivos.
          \item \cod{applicationId} definitivo (não muda depois).
          \item \cod{targetSdk} atendendo à exigência vigente.
          \item Textos em \arq{strings.xml}; tema claro e escuro revisados.
          \item Mensagens de erro, estados vazios e modo offline testados.
        \end{itemize}
      \end{conceito}
    \end{column}
    \begin{column}{0.48\textwidth}
      \begin{conceito}[Técnico]
        \begin{itemize}
          \item Build release testado em aparelhos reais (versões antigas e novas).
          \item Logs de depuração e URLs de teste removidos.
          \item Nenhum segredo no código ou no repositório.
          \item Crashlytics ativo e \arq{mapping.txt} enviado.
          \item Apenas as permissões realmente usadas.
        \end{itemize}
      \end{conceito}
    \end{column}
  \end{columns}
  \begin{atencao}[Nível de API alvo]
    Desde \textbf{31/08/2026}, apps novos e atualizações na Google Play precisam de \cod{targetSdk} \textbf{36} (Android 16) ou superior. A exigência sobe todo ano, por volta de agosto.
  \end{atencao}
\end{frame}

% =============================================================================
\section{Assinatura do app}

\begin{frame}[fragile]{Criando a chave e assinando}
  \begin{columns}[T]
    \begin{column}{0.54\textwidth}
\begin{terminal}[title=Gerando a chave de upload]
keytool -genkeypair -v -keystore upload.jks \
  -alias upload -keyalg RSA -keysize 2048 \
  -validity 10000
\end{terminal}
\begin{codigo}[title=app/build.gradle.kts]{Kotlin}
signingConfigs {
  create("release") {
    storeFile = file(System.getenv("KEYSTORE_PATH")
      ?: "upload.jks")
    storePassword = System.getenv("KEYSTORE_PASS")
    keyAlias = "upload"
    keyPassword = System.getenv("KEY_PASS")
  }
}
\end{codigo}
\begin{terminal}[title=Gerando o App Bundle]
./gradlew bundleRelease
# app/build/outputs/bundle/release/app-release.aab
\end{terminal}
    \end{column}
    \begin{column}{0.44\textwidth}
      \begin{erro}[Nunca no Git]
        O \arq{.jks} e as senhas \textbf{não} vão para o repositório. Use variáveis de ambiente, \arq{\textasciitilde/.gradle/gradle.properties} ou \emph{secrets} da CI.
      \end{erro}
      \begin{dica}[Backup]
        Guarde o \arq{.jks} e as senhas num gerenciador de senhas e em uma cópia offline.
      \end{dica}
      \begin{dica}
        Pelo Android Studio: \textbf{Build $\rightarrow$ Generate Signed App Bundle / APK} faz o mesmo com um assistente.
      \end{dica}
    \end{column}
  \end{columns}
\end{frame}

% =============================================================================
\section{Play Console}

\begin{frame}{A conta de desenvolvedor}
  \begin{columns}[T]
    \begin{column}{0.5\textwidth}
      \scriptsize
      \begin{tabular}{@{}p{1.9cm} p{2.0cm} p{2.3cm}@{}}
        \toprule
        & \cabecalho{Pessoal} & \cabecalho{Organização} \\
        \midrule
        Para & pessoa física & empresas, instituições \\
        Exige & verificação de identidade & D-U-N-S e verificação \\
        Teste fechado obrigatório & sim (contas novas) & não \\
        Nome exibido & seu nome & nome da organização \\
        \bottomrule
      \end{tabular}
      \vspace{3pt}
      \begin{conceito}[Custo]
        Taxa de registro \textbf{única}, paga ao criar a conta em \url{https://play.google.com/console}.
      \end{conceito}
    \end{column}
    \begin{column}{0.48\textwidth}
      \begin{dica}[Para projetos institucionais]
        Apps que representam o IFMA, um laboratório ou um projeto de extensão ficam melhor em uma \textbf{conta de organização} da instituição, que garante continuidade quando os alunos se formam.
      \end{dica}
      \begin{atencao}
        A conta exige verificação de identidade, telefone e e-mail de contato públicos. Mantenha esses dados atualizados: contas inativas ou com dados inválidos podem ser encerradas.
      \end{atencao}
    \end{column}
  \end{columns}
\end{frame}

\begin{frame}{Ficha da loja e declarações obrigatórias}
  \begin{columns}[T]
    \begin{column}{0.48\textwidth}
      \begin{conceito}[Ficha da loja]
        \begin{itemize}
          \item \textbf{Nome} do app (até 30 caracteres).
          \item \textbf{Descrição curta} (até 80) e \textbf{completa} (até 4000).
          \item \textbf{Ícone} 512$\times$512\,px.
          \item \textbf{Gráfico de recursos} 1024$\times$500\,px.
          \item \textbf{Capturas de tela} (celular e tablet).
          \item Categoria e e-mail de contato.
        \end{itemize}
      \end{conceito}
    \end{column}
    \begin{column}{0.48\textwidth}
      \begin{conceito}[Conteúdo do app (Política)]
        \begin{itemize}
          \item URL da \textbf{política de privacidade}.
          \item Formulário de \textbf{segurança dos dados}.
          \item \textbf{Classificação de conteúdo} (questionário IARC).
          \item \textbf{Público-alvo} e se contém anúncios.
          \item Declaração de permissões sensíveis.
          \item Acesso para o revisor (se houver login).
        \end{itemize}
      \end{conceito}
    \end{column}
  \end{columns}
  \begin{atencao}
    Se o app permite criar conta, ele deve oferecer \textbf{exclusão da conta} dentro do app e por um link na web.
  \end{atencao}
\end{frame}

% =============================================================================
\section{Faixas de teste}

\begin{frame}{Do teste interno à produção}
  \begin{center}
  \begin{tikzpicture}[node distance=6mm]
    \node[caixa,minimum width=2.6cm,minimum height=12mm] (i) {\textbf{Interno}\\ \tiny até 100 testadores\\ disponível em minutos};
    \node[caixaambar,minimum width=2.6cm,minimum height=12mm,right=of i] (f) {\textbf{Fechado}\\ \tiny lista de e-mails ou\\ Grupo Google};
    \node[caixapetroleo,minimum width=2.6cm,minimum height=12mm,right=of f] (a) {\textbf{Aberto}\\ \tiny qualquer pessoa pode\\ entrar no teste};
    \node[caixaescura,minimum width=2.6cm,minimum height=12mm,right=of a] (p) {\textbf{Produção}\\ \tiny todos os usuários};
    \draw[seta] (i) -- (f); \draw[seta] (f) -- (a); \draw[seta] (a) -- (p);
    \draw[setaambar,rounded corners] (f.south) -- ++(0,-5mm) -| node[rotulo,pos=0.25,below]{contas pessoais novas: 12 testadores por 14 dias antes de pedir acesso à produção} (p.south);
  \end{tikzpicture}
  \end{center}
  \begin{columns}[T]
    \begin{column}{0.48\textwidth}
      \begin{atencao}[Regra dos 12 testadores]
        Contas \textbf{pessoais} criadas após 13/11/2023 precisam manter um teste fechado com \textbf{pelo menos 12 testadores inscritos por 14 dias seguidos} antes de solicitar acesso à produção.
      \end{atencao}
    \end{column}
    \begin{column}{0.48\textwidth}
      \begin{dica}[Na prática]
        \begin{itemize}
          \item Recrute 15--20 testadores: se alguém sair, a contagem pode recomeçar.
          \item Uma turma do curso é um ótimo grupo de testadores!
          \item Publique atualizações durante o teste e registre o feedback.
        \end{itemize}
      \end{dica}
    \end{column}
  \end{columns}
\end{frame}

% =============================================================================
\section{Lançamento e operação}

\begin{frame}{Na Play: lançamento gradual e operação}
  \begin{columns}[T]
    \begin{column}{0.48\textwidth}
      \begin{conceito}[Lançamento gradual na prática]
        Ao promover a versão para produção, escolha a porcentagem inicial e aumente-a aos poucos, acompanhando o Android vitals. Problema? Interrompa o lançamento: quem ainda não atualizou continua na versão anterior.
      \end{conceito}
      \begin{atencao}[Não existe ``voltar versão'']
        Quem já atualizou fica com a versão nova. A correção é sempre outra versão, com \cod{versionCode} maior.
      \end{atencao}
    \end{column}
    \begin{column}{0.48\textwidth}
      \begin{dica}[Depois de publicado]
        \begin{itemize}
          \item Responda às avaliações: a resposta é pública.
          \item Use \textbf{In-App Updates} para pedir a atualização de versões críticas.
          \item Use \textbf{In-App Review} para pedir avaliação no momento certo.
          \item Leia as notas de política da Play: as regras mudam.
          \item Planeje a atualização anual do \cod{targetSdk}.
        \end{itemize}
      \end{dica}
    \end{column}
  \end{columns}
\end{frame}

\begin{frame}{Distribuição fora da Play e verificação de desenvolvedor}
  \begin{columns}[T]
    \begin{column}{0.52\textwidth}
      \begin{atencao}[Novidade de 2026 no Brasil]
        Desde \textbf{30/09/2026}, aparelhos Android certificados no Brasil (e também na Indonésia, em Singapura e na Tailândia) só instalam normalmente apps \textbf{registrados por desenvolvedores verificados}, inclusive quando distribuídos fora da Google Play (APK em site, lojas alternativas). A expansão global está prevista para 2027.
      \end{atencao}
      \begin{conceito}[Como se registrar]
        Pela \textbf{Play Console} (apps na loja) ou pela \textbf{Android Developer Console} (apps só fora da loja), verificando a identidade e registrando o nome do pacote e a chave de assinatura.
      \end{conceito}
    \end{column}
    \begin{column}{0.46\textwidth}
      \begin{dica}[Estudantes e hobbistas]
        Existem contas de \textbf{distribuição limitada}, pensadas para estudantes e projetos pessoais que instalam apps em poucos aparelhos.
      \end{dica}
      \begin{conceito}[O que não muda para o desenvolvimento]
        Instalar pelo Android Studio ou via \textbf{ADB} continua funcionando. Usuários avançados também têm um fluxo especial para instalar apps não verificados.
      \end{conceito}
      \begin{dica}
        Para distribuir o LabMonitor só no laboratório via APK, registre o app antes.
      \end{dica}
    \end{column}
  \end{columns}
\end{frame}

% =============================================================================
\section{Automatizando com CI/CD}

\begin{frame}[fragile]{Publicação automática}
  \begin{columns}[T]
    \begin{column}{0.56\textwidth}
\begin{codigo}[title=.github/workflows/release.yml (trecho)]{bash}
on:
  push:
    tags: ["v*"]          # ex.: git tag v1.2.0
jobs:
  release:
    runs-on: ubuntu-latest
    steps:
      - uses: actions/checkout@v4
      - uses: actions/setup-java@v4
        with: { distribution: temurin, java-version: 17 }
      - name: Restaurar keystore
        run: echo "$KEYSTORE_B64" | base64 -d > upload.jks
        env: { KEYSTORE_B64: "${{ secrets.KEYSTORE_B64 }}" }
      - name: Gerar AAB
        run: ./gradlew bundleRelease
        env:
          KEYSTORE_PASS: ${{ secrets.KEYSTORE_PASS }}
          KEY_PASS: ${{ secrets.KEY_PASS }}
      # enviar à faixa interna (Gradle Play Publisher
      # ou Fastlane supply) com uma conta de serviço
\end{codigo}
    \end{column}
    \begin{column}{0.42\textwidth}
      \begin{conceito}[Pipeline típico]
        \emph{push} de uma tag $\rightarrow$ testes $\rightarrow$ build assinado $\rightarrow$ envio à faixa interna $\rightarrow$ promoção manual para produção.
      \end{conceito}
      \begin{dica}[Ferramentas]
        \textbf{Gradle Play Publisher} (plugin Gradle) ou \textbf{Fastlane} (\cod{supply}) usam a API da Play Console com uma \emph{conta de serviço} do Google Cloud.
      \end{dica}
      \begin{atencao}
        Segredos ficam nos \emph{secrets} do repositório, nunca no YAML.
      \end{atencao}
    \end{column}
  \end{columns}
\end{frame}

% =============================================================================
\section{Fechamento}

\begin{frame}{Resumo}
  \begin{resumo}
    \begin{itemize}
      \item Release = R8, \cod{versionCode} crescente, \cod{targetSdk} atualizado e testes no build final.
      \item Play App Signing: você guarda a chave de upload; o Google guarda a chave do app.
      \item A loja exige ficha completa, política de privacidade, segurança dos dados e classificação.
      \item Faixas interna, fechada e aberta antes da produção; contas pessoais novas precisam de 12 testadores por 14 dias.
      \item No Brasil, desde 30/09/2026, apps precisam ser de desenvolvedores verificados mesmo fora da Play.
      \item LGPD: mínimo de dados, transparência, consentimento e exclusão. CI/CD automatiza tudo.
    \end{itemize}
  \end{resumo}
\end{frame}

\begin{frame}{Glossário (1/2)}
  \begin{glossario}
    Build de release & versão otimizada e assinada do app, destinada ao usuário final \\
    R8 & encolhe, otimiza e ofusca o código no build de release \\
    Assinatura digital & \emph{hash} do pacote cifrado com a chave privada; prova integridade e autoria \\
    \emph{Keystore} & arquivo protegido por senha que guarda as chaves de assinatura \\
    Play App Signing & o Google guarda a chave do app; você assina os envios com a chave de upload \\
    AAB & \emph{Android App Bundle}: formato de envio do qual a Play gera APKs por aparelho \\
    \cod{versionCode} & número interno que precisa crescer a cada versão enviada \\
    \arq{mapping.txt} & traduz pilhas ofuscadas pelo R8 de volta para nomes legíveis \\
    SemVer & versionamento semântico do \cod{versionName}: MAJOR.MINOR.PATCH \\
  \end{glossario}
\end{frame}

\begin{frame}{Glossário (2/2)}
  \begin{glossario}[3.9cm]
    Faixa de teste & canal da Play para distribuir a grupos: interna, fechada ou aberta \\
    \emph{Data safety} & declaração dos dados coletados, exibida na ficha da loja \\
    Lançamento gradual & liberação da versão para uma porcentagem crescente de usuários \\
    Android vitals & falhas, ANRs e desempenho medidos nos aparelhos dos usuários \\
    Verificação de desenvolvedor & registro de identidade exigido para instalar apps em aparelhos certificados \\
    ASO & \emph{App Store Optimization}: ser encontrado e escolhido na loja \\
    CI/CD & compilar, testar e publicar automaticamente a cada mudança \\
  \end{glossario}
\end{frame}

\begin{frame}{Questões de revisão}
  \begin{questoes}
    \item O que a assinatura digital garante, e por que perder a chave pode impedir atualizações?
    \item Qual a diferença entre a chave de upload e a chave do app no Play App Signing?
    \item Por que o \cod{versionCode} precisa sempre aumentar? Como ``desfazer'' uma versão com defeito?
    \item Em que o AAB difere do APK, e que vantagem ele traz para o usuário?
    \item Por que é preciso guardar o \arq{mapping.txt} de cada versão?
    \item Uma conta pessoal nova quer publicar. Que etapa de teste a Play exige antes da produção?
    \item O que mudou, desde 30/09/2026, para distribuir um APK fora da Play no Brasil?
    \item Como um lançamento gradual reduz o risco de uma versão nova?
  \end{questoes}
\end{frame}

\begin{frame}{Exercícios}
  \begin{exercicio}[Exercício 1 --- Release]
    Gere o AAB assinado do LabMonitor, instale o build release com \cod{bundletool} e verifique se tudo funciona com o R8 ativo.
  \end{exercicio}
  \begin{exercicio}[Exercício 2 --- Política]
    Escreva a política de privacidade do LabMonitor e preencha um rascunho do formulário de segurança dos dados com base no que o app realmente coleta.
  \end{exercicio}
  \begin{exercicio}[Exercício 3 --- Ficha]
    Produza ícone, gráfico de recursos, 4 capturas de tela e as descrições curta e completa do app.
  \end{exercicio}
  \begin{projeto}[LabMonitor --- Etapa 15]
    Publique o app na faixa de \textbf{teste fechado} com a turma como testadores, colete o feedback por 14 dias e lance a versão 1.0.
  \end{projeto}
\end{frame}

\begin{frame}{Para aprofundar}
  \small
  \begin{itemize}
    \item Preparar o app para lançamento: \url{https://developer.android.com/studio/publish/preparing}
    \item Assinatura de apps: \url{https://developer.android.com/studio/publish/app-signing}
    \item Requisitos de nível de API: \url{https://developer.android.com/google/play/requirements/target-sdk}
    \item Testes para contas pessoais novas: \url{https://support.google.com/googleplay/android-developer/answer/14151465}
    \item Verificação de desenvolvedor: \url{https://developer.android.com/developer-verification}
    \item LGPD (texto da lei): \url{https://www.planalto.gov.br/ccivil_03/_ato2015-2018/2018/lei/l13709.htm}
  \end{itemize}
\end{frame}

\begin{frame}{Respostas das questões de revisão}
  \begin{respostas}
    \item Integridade (o pacote não foi alterado) e autenticidade (vem do mesmo autor). O sistema só aceita atualizações assinadas pela mesma chave; sem ela, não há como provar a autoria.
    \item A de upload assina o que você envia e pode ser trocada; a do app, guardada pelo Google, assina os APKs entregues aos aparelhos.
    \item A loja e o sistema usam o \cod{versionCode} para saber qual build é mais novo e recusam números menores. Para desfazer, publica-se uma versão nova, com número maior, contendo a correção.
    \item O APK contém tudo para todos os aparelhos; do AAB a Play gera APKs só com o código, as imagens e os idiomas de cada aparelho: download menor.
    \item O R8 renomeia classes e métodos; sem o mapeamento, as pilhas de falha no Android vitals e no Crashlytics ficam ilegíveis.
    \item Um teste fechado com pelo menos 12 testadores inscritos por 14 dias seguidos, exigido de contas pessoais criadas depois de 13/11/2023.
    \item Aparelhos certificados só instalam normalmente apps de desenvolvedores verificados, mesmo fora da Play: é preciso verificar a identidade e registrar o pacote e a chave.
    \item Um defeito atinge só a fração inicial de usuários; os indicadores de saúde decidem se o lançamento avança ou é interrompido.
  \end{respostas}
\end{frame}

\encerramento{Módulo 16 (comum) --- Backend terceirizado: Firebase, Supabase, AWS Amplify, Appwrite e outros BaaS do mercado.}

\end{document}
